\NormalFont\ShortTitle{1 João}
{\MT Primeira Carta Universal de João

\par }\ChapOne{1}{\PP \VerseOne{1}O que era desde o princípio, o que ouvimos, o que vimos com os nossos olhos, o que temos contemplado, e nossas mãos tocaram, quanto à Palavra da vida,
\VS{2}(pois a vida já foi manifesta, e nós a vimos, demos testemunho, e vos anunciamos a vida eterna, que estava com o Pai, e a nós foi manifesta)
\VS{3}o que vimos e ouvimos, isso vos anunciamos, para que também vós tenhais comunhão conosco. E a nossa comunhão
{\ADD{é}} com o Pai, e com o seu Filho Jesus Cristo.
\VS{4}Estas coisas vos escrevemos para que a vossa alegria seja completa.
\VS{5}E esta é a mensagem que dele ouvimos, e vos anunciamos: que Deus é luz, e não há nele nada de trevas.
\VS{6}Se dissermos que temos comunhão com ele, e andarmos em trevas, mentimos, e não praticamos a verdade.
\VS{7}Porém, se andamos na luz, como ele está na luz, temos comunhão uns com os outros, e o sangue de Jesus Cristo, seu Filho, nos purifica de todo pecado.
\VS{8}Se dissermos que não temos pecado, enganamos a nós mesmos, e a verdade não está em nós.
\VS{9}Se confessarmos nossos pecados, ele é fiel e justo para nos perdoar os pecados e nos purificar de toda injustiça.
\VS{10}Se dissermos que não pecamos, nós o fazemos de mentiroso, e a palavra dele não está em nós.

\par }\Chap{2}{\PP \VerseOne{1}Meus filhinhos, vos escrevo estas coisas para que não pequeis; mas se alguém pecar, temos um advogado com o Pai: Jesus Cristo, o justo.
\VS{2}E ele é o sacrifício para perdão dos nossos pecados; e não somente dos nossos, mas também dos de todo o mundo.
\VS{3}E nisto sabemos que o conhecemos, se guardamos os seus mandamentos.
\VS{4}Aquele que diz: “Eu o conheço”, mas não guarda os seus mandamentos, é mentiroso, e nele não está a verdade.
\VS{5}Mas quem guarda a sua palavra, nele verdadeiramente o amor de Deus tem se aperfeiçoado; nisso sabemos que estamos nele.
\VS{6}Aquele que diz que está nele também deve andar como ele andou.
\VS{7}Irmãos, não vos escrevo um mandamento novo, mas sim o mandamento antigo, que tivestes desde o princípio. Este mandamento antigo é a palavra que ouvistes desde o princípio.
\VS{8}Mas também vos escrevo um mandamento novo, que é verdadeiro nele e em vós; pois as trevas passaram, e a verdadeira luz já ilumina.
\VS{9}Aquele que diz que está na luz, e odeia seu irmão, até agora está nas trevas.
\VS{10}Aquele que ama o seu irmão está na luz, e nele não há motivo algum de tropeço.
\VS{11}Mas aquele que odeia seu irmão está nas trevas, anda nas trevas, e não sabe para onde vai, pois as trevas lhe cegaram os olhos.
\VS{12}Eu vos escrevo, filhinhos, porque pelo nome dele os vossos pecados são perdoados.
\VS{13}Eu vos escrevo, pais, porque conheceis aquele que é desde o princípio. Eu vos escrevo, jovens, porque vencestes o maligno. Eu vos escrevo, filhinhos, porque conheceis o Pai.
\VS{14}Eu vos escrevi, pais, porque conheceis aquele que é desde o princípio. Eu vos escrevi, jovens, porque sois fortes, a palavra de Deus está em vós, e vencestes o maligno.
\VS{15}Não ameis o mundo, nem as coisas que há no mundo. Se alguém ama o mundo, o amor do Pai não está nele.
\VS{16}Pois tudo o que há no mundo, a cobiça da carne, a cobiça dos olhos, e a ostentação da vida, não é do Pai, mas do mundo.
\VS{17}O mundo passa, e também a sua cobiça; mas aquele que faz a vontade de Deus permanece para sempre.
\VS{18}Filhinhos, é a última hora. E como ouvistes que o anticristo vem, já agora muitos anticristos chegaram. Por isso sabemos que é a última hora.
\VS{19}Eles saíram de nós, mas não eram dos nossos, pois se fossem dos nossos, continuariam conosco; mas
{\ADD{o que aconteceu}} foi para tornar claro que todos eles não eram dos nossos.
\VS{20}Vós, porém, tendes a unção do Santo, e conheceis todas as coisas.
\VS{21}Eu vos escrevi, não porque não conheceis a verdade, mas sim, porque a conheceis, e porque nenhuma mentira vem da verdade.
\VS{22}Quem é o mentiroso, senão aquele que nega que Jesus é o Cristo? Esse é o anticristo, que nega o Pai e o Filho.
\VS{23}Quem nega o Filho também não tem o Pai; aquele que confessa o filho também tem o Pai.
\VS{24}Portanto o que ouvistes desde o princípio permaneça em vós. Se o que ouvistes desde o princípio permanecer em vós, também permanecereis no Filho e no Pai.
\VS{25}E esta é a promessa que ele nos prometeu: a vida eterna.
\VS{26}Estas coisas vos escrevi acerca dos que vos tentam vos enganar.
\VS{27}E a unção que recebeste dele continua em vós, e não tendes necessidade de que alguém vos ensine; mas, como a mesma unção vos ensina tudo, e é verdadeira, e não mentira, e como ela vos ensinou,
{\ADD{assim}} nele permanecereis.
\VS{28}E agora, filhinhos, permanecei nele; a fim de que, quando ele se manifestar, tenhamos confiança, e não sejamos envergonhados diante dele na sua vinda.
\VS{29}Se sabeis que ele é justo, sabeis que todo aquele que pratica a justiça é nascido dele.

\par }\Chap{3}{\PP \VerseOne{1}Vede quão grande amor o Pai tem nos dado, que fôssemos chamados filhos de Deus. Por isso o mundo não nos conhece, pois não conhece a ele.
\VS{2}Amados, agora somos filhos de Deus, e ainda não está manifesto o que seremos. Porém sabemos que, quando ele se manifestar, seremos semelhantes a ele, pois o veremos como ele é.
\VS{3}E todo aquele que tem nele essa esperança purifica a si mesmo, como também ele é puro.
\VS{4}Todo aquele que pratica o pecado também pratica injustiça, pois o pecado é injustiça.
\VS{5}E sabeis que ele apareceu para tirar os nossos pecados; e nele não há pecado.
\VS{6}Todo aquele que nele permanece não pratica o pecado; todo aquele que costuma pecar não o viu nem o conheceu.
\VS{7}Filhinhos, ninguém vos engane. Quem pratica a justiça é justo, assim como ele é justo.
\VS{8}Quem pratica o pecado é do diabo, pois o diabo peca desde o princípio. Para isto o Filho de Deus se manifestou, para desfazer as obras do diabo.
\VS{9}Todo aquele que é nascido de Deus não pratica o pecado, pois a sua semente reside nele; e não pode praticar o pecado, porque é nascido de Deus.
\VS{10}Nisto são reconhecíveis os filhos de Deus, e os filhos do diabo: todo aquele que não pratica a justiça, e não ama o seu irmão, não é de Deus.
\VS{11}Pois esta é a mensagem que ouvistes desde o princípio: que nos amemos uns aos outros.
\VS{12}Não
{\ADD{sejamos}} como Caim,
{\ADD{ que}} era do maligno, e matou o seu irmão. E por que o matou? Porque as suas obras eram más, e as do seu irmão justas.
\VS{13}Meus irmãos, não vos surpreendeis se o mundo vos odeia.
\VS{14}Nós sabemos que já passamos da morte para a vida, pois amamos os irmãos. Quem não ama{\ADD{seu}} irmão permanece na morte.
\VS{15}Todo aquele que odeia o seu irmão é homicida. E sabeis que nenhum homicida tem a vida eterna habitando nele.
\VS{16}Nisto conhecemos o amor: ele deu a sua vida por nós. E nós devemos dar as
{\ADD{nossas}} vidas pelos irmãos.
\VS{17}Se alguém tiver bens do mundo, e vir o seu irmão em necessidade, e não se compadecer dele, como pode o amor de Deus estar nele?
\VS{18}Meus filhinhos, amemos não de palavra, nem de língua, mas sim com ação e verdade.
\VS{19}E nisto sabemos que somos da verdade, e teremos segurança nos nossos corações diante dele.
\VS{20}Pois, se o nosso coração
{\ADD{nos}} condena, maior é Deus que o nosso coração, e conhece tudo.
\VS{21}Amados, se nosso coração não nos condena, temos confiança diante de Deus.
\VS{22}E qualquer coisa que pedirmos, dele receberemos; porque guardamos os seus mandamentos, e fazemos o que lhe agrada.
\VS{23}E este é o seu mandamento: que creiamos no nome do seu Filho Jesus Cristo, e nos amemos uns aos outros, como ele nos mandou.
\VS{24}E a pessoa que guarda os seus mandamentos está nele, e ele nela. E nisto sabemos que ele está em nós: pelo Espírito que ele nos deu.

\par }\Chap{4}{\PP \VerseOne{1}Amados, não creiais em todo espírito, mas provai se os espíritos são da parte de Deus; porque muitos falsos profetas têm vindo ao mundo.
\VS{2}Nisto conhecereis o Espírito de Deus: todo espírito que confessa que Jesus Cristo veio em carne é de Deus;
\VS{3}e todo espírito que não confessa que Jesus Cristo veio em carne não é de Deus; e este é o do anticristo, do qual já ouvistes que virá, e que já agora está no mundo.
\VS{4}Filhinhos, vós sois de Deus, e tendes vencido esses; pois maior é o que está em vós do que o que está no mundo.
\VS{5}Eles são do mundo, por isso falam do mundo, e o mundo os ouve.
\VS{6}Nós somos de Deus. Aquele que conhece Deus ouve-nos; aquele que não é de Deus não nos ouve. Nisto conhecemos o Espírito da verdade e o espírito do erro.
\FTNT{*}{{\FR 4:6 }
Ou: engano
}
\VS{7}Amados, nos amemos uns aos outros; porque o amor é de Deus, e todo aquele que ama é nascido de Deus e conhece a Deus.
\VS{8}Aquele que não ama não conhece a Deus, porque Deus é amor.
\VS{9}Nisto se manifestou o amor de Deus por nós:
\FTNT{†}{{\FR 4:9 }
Ou: entre nós
} que Deus enviou o seu Filho unigênito ao mundo, para que por meio dele vivamos.
\VS{10}Nisto está o amor, não que nós tenhamos amado a Deus, mas sim, que ele nos amou, e enviou seu Filho como sacrifício para perdão dos nossos pecados.
\VS{11}Amados, se Deus assim nos amou, também nós devemos amar uns aos outros.
\VS{12}Ninguém jamais viu Deus; se nos amamos uns aos outros, Deus está em nós, e o amor dele em nós é aperfeiçoado.
\VS{13}Nisto sabemos que estamos nele, e ele em nós, porque ele nos deu do seu Espírito.
\VS{14}E vimos e damos testemunho de que o Pai enviou o Filho
{\ADD{para ser}} Salvador do mundo.
\VS{15}Todo aquele que confessar que Jesus é o Filho de Deus, Deus está nele, e ele em Deus.
\VS{16}E nós conhecemos e cremos no amor que Deus tem por nós. Deus é amor, e quem está no amor está em Deus, e Deus nele.
\VS{17}Nisto o amor é aperfeiçoado em nós, para que no dia do juízo tenhamos confiança, pois, como ele é, assim somos nós neste mundo.
\VS{18}No amor não há medo; pelo contrário, o perfeito amor expulsa o medo; pois o medo está relacionado ao castigo; e o que teme não está aperfeiçoado em amor.
\VS{19}Nós o amamos, porque ele nos amou primeiro.
\VS{20}Se alguém diz: “Eu amo a Deus”, e odeia seu irmão, é mentiroso. Pois quem não ama seu irmão, ao qual viu, como pode amar a Deus, a quem nunca viu?
\VS{21}E dele temos este mandamento: quem ama a Deus, ame também seu irmão.

\par }\Chap{5}{\PP \VerseOne{1}Todo aquele que crê que Jesus é o Cristo é nascido de Deus; e todo aquele que ama ao que gerou, também ama ao que dele é nascido.
\VS{2}Nisto sabemos que amamos os filhos de Deus: quando amamos a Deus, e guardamos os seus mandamentos.
\VS{3}Pois este é o amor de Deus: que guardemos os seus mandamentos. E os seus mandamentos não são pesados.
\VS{4}Pois todo aquele que é nascido de Deus vence o mundo; e esta é a vitória que vence o mundo: a nossa fé.
\VS{5}Quem é o que vence o mundo, senão aquele que crê que Jesus é o Filho de Deus?
\VS{6}Este é aquele que veio por água e sangue: Jesus, o Cristo. Não só pela água, mas pela água e
{\ADD{pelo}} sangue. E o Espírito é o que lhe dá testemunho, porque o Espírito é a verdade.
\VS{7}Pois são três os que dão testemunho no céu: o Pai, a Palavra, e o Espírito Santo; e estes três são um.
\VS{8}E três são os que testemunham na terra: o Espírito, a água e o sangue; e estes três estão em unidade.
\VS{9}Se recebemos o testemunho humano, o testemunho de Deus é maior; pois este é o testemunho de Deus, que testificou de seu Filho.
\VS{10}Quem crê no Filho de Deus tem testemunho em si mesmo; quem não crê em Deus o fez de mentiroso, porque não creu no testemunho que Deus deu acerca do seu Filho.
\VS{11}E este é o testemunho: que Deus nos deu a vida eterna; e essa vida está no seu Filho.
\VS{12}Quem tem o Filho tem a vida; quem não tem o Filho de Deus não tem a vida.
\VS{13}Estas coisas escrevi para vós que credes no nome do Filho de Deus, para que saibais que tendes a vida eterna, e que creiais no nome do Filho de Deus.
\VS{14}E esta é a confiança que diante dele: que se pedirmos alguma coisa conforme a sua vontade, ele nos ouve.
\VS{15}E se sabemos que nos ouve em tudo o que pedimos, sabemos que alcançamos as coisas que lhe pedimos.
\VS{16}Se alguém vê o seu irmão cometer um pecado que não é para a morte, orará, e lhe dará vida, para os que não pecarem para a morte. Há pecado para a morte, pelo qual não digo que ore.
\VS{17}Toda injustiça é pecado; e há pecado que não é para a morte.
\VS{18}Sabemos que todo aquele que é nascido de Deus não pratica o pecado; mas o que é gerado de Deus guarda a si mesmo, e o maligno não o toca.
\VS{19}Sabemos que somos de Deus, e que o mundo todo jaz no maligno.
\VS{20}E sabemos que o Filho de Deus veio, e nos deu entendimento para conhecermos o Verdadeiro; e no Verdadeiro estamos, em seu Filho Jesus Cristo. Este é o verdadeiro Deus, e a vida eterna.
\VS{21}Filhinhos, guardai-vos dos ídolos. Amém.\par }